\documentclass[14pt,a4paper]{article}

% --- Подключение общего стиля из родительской папки ---
\makeatletter
\def\input@path{{../}}
\makeatother
\usepackage{common}

% --- Папка для картинок этого документа ---
\graphicspath{{fig/}}

\begin{document}

% ===================== ТИТУЛЬНЫЙ ЛИСТ =====================
\begin{titlepage}
\begin{center}
МИНИСТЕРСТВО НАУКИ И ВЫСШЕГО ОБРАЗОВАНИЯ РОССИЙСКОЙ ФЕДЕРАЦИИ\\
ФЕДЕРАЛЬНОЕ ГОСУДАРСТВЕННОЕ АВТОНОМНОЕ ОБРАЗОВАТЕЛЬНОЕ УЧРЕЖДЕНИЕ\\
ВЫСШЕГО ОБРАЗОВАНИЯ\\
«САМАРСКИЙ НАЦИОНАЛЬНЫЙ ИССЛЕДОВАТЕЛЬСКИЙ УНИВЕРСИТЕТ\\
ИМЕНИ АКАДЕМИКА С. П. КОРОЛЕВА»\\
(Самарский университет)\\[12pt]
Кафедра \underline{\hspace{6cm}}
\end{center}

\vfill

\begin{center}
ПОЯСНИТЕЛЬНАЯ ЗАПИСКА\\
к курсовому проекту по дисциплине\\
«Технологии программирования»\\[12pt]
\textbf{Интерпретатор BASIC-подобного языка}\\
Вариант 3-2-1-A
\end{center}

\vfill

\begin{flushright}
\begin{tabular}{l}
Студент группы \underline{\hspace{2.5cm}} \\
\underline{\hspace{5cm}} \\
Руководитель \underline{\hspace{2.5cm}} \\
\underline{\hspace{5cm}} \\
Оценка \underline{\hspace{3cm}} \\
Дата \underline{\hspace{3cm}}
\end{tabular}
\end{flushright}

\vfill

\begin{center}
Самара 2024
\end{center}
\end{titlepage}

% ===================== ЗАДАНИЕ =====================
\structtitle{Задание}

% ФОТО: скан/фото выданного задания.
% Положите файл в zapiska/fig/zadanie.png и раскомментируйте:
% \begin{figure}[H]
%   \centering
%   \includegraphics[width=0.9\textwidth]{zadanie.png}
% \end{figure}

\begin{figure}[H]
  \centering
  \fbox{\parbox[c][9cm][c]{0.85\textwidth}{\centering
    \textbf{[МЕСТО ДЛЯ ФОТО]}\\[4pt]
    Скан выданного задания на курсовой проект}}
\end{figure}

% ===================== РЕФЕРАТ =====================
\structtitle{Реферат}

Пояснительная записка содержит \underline{\hspace{1cm}}~с., \underline{\hspace{1cm}}~рис.,
\underline{\hspace{1cm}}~табл., \underline{\hspace{1cm}}~источников, \underline{\hspace{1cm}}~прил.

\vspace{6pt}
КЛЮЧЕВЫЕ СЛОВА: ИНТЕРПРЕТАТОР, BASIC-ПОДОБНЫЙ ЯЗЫК, ЛЕКСИЧЕСКИЙ
АНАЛИЗАТОР, КОНЕЧНЫЙ АВТОМАТ, СИНТАКСИЧЕСКИЙ АНАЛИЗАТОР,
РЕКУРСИВНЫЙ СПУСК, ПОЛИЗ, СТЕКОВАЯ МАШИНА, ХЕШ-ТАБЛИЦА,
МЕТОД ЦЕПОЧЕК, СТРОГАЯ ТИПИЗАЦИЯ, ОТЛАДЧИК.

\vspace{6pt}
Объект разработки — интерпретатор упрощённого императивного
языка программирования с BASIC-подобным синтаксисом (вариант 3-2-1-A:
цикл с постусловием, хеш-таблица с разрешением коллизий методом
цепочек, строгая типизация).

Цель работы — разработка интерпретатора с возможностью
редактирования и отладки программ, обеспечивающего лексический
и синтаксический анализ, построение внутренней формы представления
программы в виде ПОЛИЗа, выполнение программы на стековой машине,
а также просмотр текущей инструкции, стека и значений переменных.

В результате работы реализованы: лексический анализатор на основе
конечного автомата; синтаксический анализатор методом рекурсивного
спуска с семантическими процедурами генерации ПОЛИЗа; таблица
идентификаторов в виде хеш-таблицы с цепочками; стековая машина;
отладчик с пошаговым выполнением и точками останова; графический
интерфейс.

Область применения — учебные курсы по технологиям программирования,
теории компиляторов и языкам программирования.

% ===================== СОДЕРЖАНИЕ =====================
\structtitle{Содержание}

\tableofcontents

% ===================== ВВЕДЕНИЕ =====================
\structtitle{Введение}

В соответствии с заданием на курсовой проект требуется разработать
интерпретатор упрощённого императивного языка с возможностью
редактирования и отладки программ. Вариант задания — 3-2-1-A:
BASIC-подобный синтаксис; оператор цикла с постусловием; таблица
идентификаторов на основе хеш-таблицы с разрешением коллизий
методом цепочек; строгая типизация.

Актуальность работы определяется тем, что построение интерпретатора
позволяет на практике освоить все этапы трансляции: лексический
анализ, синтаксический анализ, генерацию внутреннего представления
и его выполнение. Эти этапы лежат в основе современных компиляторов
и интерпретаторов, а также инструментов статического анализа.

Цель работы — разработка интерпретатора, удовлетворяющего
требованиям задания на оценку «5»: поддержка функций, оператора
множественного выбора, безусловного перехода, графического
интерфейса, пошагового выполнения, точек останова, редактирования
переменных и синхронизации ПОЛИЗа с исходным текстом.

Для достижения цели решаются следующие задачи:
\begin{itemize}
  \item формально описать язык (ключевые слова, лексика, грамматика,
        правила типизации);
  \item спроектировать архитектуру интерпретатора;
  \item реализовать лексический анализатор на основе конечного
        автомата;
  \item реализовать синтаксический анализатор методом рекурсивного
        спуска с семантическими процедурами генерации ПОЛИЗа;
  \item реализовать таблицу идентификаторов в виде хеш-таблицы
        с цепочками;
  \item реализовать стековую машину для выполнения ПОЛИЗа;
  \item реализовать отладчик и графический интерфейс;
  \item провести тестирование.
\end{itemize}

% ===================== 1 ФОРМАЛЬНОЕ ОПИСАНИЕ ЯЗЫКА =====================
\section{Формальное описание языка}

\subsection{Ключевые слова}

Язык поддерживает следующие ключевые слова (регистр не важен):
DIM, AS, INTEGER, REAL, STRING, IF, THEN, ELSE, ENDIF, DO, LOOP,
WHILE, UNTIL, FOR, TO, STEP, NEXT, SELECT, CASE, ENDSELECT, GOTO,
RETURN, FUNCTION, ENDFUNCTION, CALL, PRINT, INPUT, AND, OR, NOT.

\subsection{Лексика}

Лексические единицы языка:
\begin{itemize}
  \item идентификатор: \verb|[A-Za-z_][A-Za-z0-9_]*|;
  \item целая константа: \verb|[0-9]+|;
  \item вещественная константа: \verb|[0-9]+\.[0-9]+|;
  \item строковая константа: \verb|"..."| (экранирование \verb|""|);
  \item комментарий: \verb|'| до конца строки;
  \item операторы: \verb|+ - * / > < >= <= = <>|;
  \item разделители: \verb|( ) , :|, конец строки.
\end{itemize}

\subsection{Грамматика}

Грамматика языка в форме Бэкуса~---~Наура:
\begin{verbatim}
program        = { function | statement } EOF ;
function       = "FUNCTION" ident "(" [ param { "," param } ] ")"
                 { statement } "ENDFUNCTION" ;
param          = ident "AS" type ;
statement      = var_decl | assignment | print_stmt | input_stmt
               | if_stmt | do_stmt | for_stmt | select_stmt
               | goto_stmt | return_stmt | call_stmt | label ;
var_decl       = "DIM" ident "AS" type { "," ident "AS" type } ;
type           = "INTEGER" | "REAL" | "STRING" ;
assignment     = ident ":=" expression ;
print_stmt     = "PRINT" [ expression { "," expression } ] ;
input_stmt     = "INPUT" ident { "," ident } ;
if_stmt        = "IF" expression "THEN" { statement }
                 [ "ELSE" { statement } ] "ENDIF" ;
do_stmt        = "DO" { statement } "LOOP" ( "WHILE" | "UNTIL" ) expression ;
for_stmt       = "FOR" ident "=" expression "TO" expression
                 [ "STEP" expression ] { statement } "NEXT" ;
select_stmt    = "SELECT" "CASE" expression
                 { "CASE" case_label { "," case_label } ":" { statement } }
                 [ "CASE" "ELSE" ":" { statement } ] "ENDSELECT" ;
case_label     = const | const "TO" const ;
goto_stmt      = "GOTO" ident ;
return_stmt    = "RETURN" [ expression ] ;
call_stmt      = "CALL" ident "(" [ expression { "," expression } ] ")" ;
label          = ident ":" ;
expression     = logic_or ;
logic_or       = logic_and { "OR" logic_and } ;
logic_and      = logic_not { "AND" logic_not } ;
logic_not      = [ "NOT" ] comparison ;
comparison     = additive [ ( ">" | "<" | ">=" | "<=" | "=" | "<>" ) additive ] ;
additive       = multiplicative { ( "+" | "-" ) multiplicative } ;
multiplicative = unary { ( "*" | "/" ) unary } ;
unary          = [ "-" ] primary ;
primary        = const | ident | call_expr | "(" expression ")" ;
call_expr      = ident "(" [ expression { "," expression } ] ")" ;
\end{verbatim}

\subsection{Строгая типизация}

Правила типизации:
\begin{itemize}
  \item преобразование INTEGER $\to$ REAL разрешено (расширяющее);
  \item преобразование REAL $\to$ INTEGER запрещено;
  \item STRING не участвует в арифметических операциях;
  \item в числовую переменную нельзя присвоить STRING и наоборот.
\end{itemize}

% ===================== 2 АРХИТЕКТУРА =====================
\section{Архитектура}

Модули проекта:
\begin{itemize}
  \item \texttt{tokens.py} — типы токенов;
  \item \texttt{lexer.py} — лексический анализатор (конечный автомат);
  \item \texttt{parser.py} — синтаксический анализатор (рекурсивный
        спуск) с семантическими процедурами генерации ПОЛИЗа;
  \item \texttt{poliz.py} — операции ПОЛИЗа, дизассемблер;
  \item \texttt{symbol\_table.py} — хеш-таблица с цепочками;
  \item \texttt{interpreter.py} — стековая машина;
  \item \texttt{debugger.py} — отладчик;
  \item \texttt{main.py} — графический интерфейс (Tkinter).
\end{itemize}

% ФОТО: схема архитектуры. Положите в zapiska/fig/arch.png
\begin{figure}[H]
  \centering
  \fbox{\parbox[c][6cm][c]{0.8\textwidth}{\centering
    \textbf{[МЕСТО ДЛЯ ФОТО]}\\[4pt]
    Схема архитектуры: Lexer $\to$ Parser $\to$ POLIZ $\to$ Interpreter\\
    с ответвлениями в SymbolTable и Debugger}}
  \caption{Архитектура интерпретатора}
  \label{fig:arch}
\end{figure}

% ===================== 3 РЕАЛИЗАЦИЯ =====================
\section{Реализация}

\subsection{Лексический анализатор}

Лексический анализатор реализован как детерминированный конечный
автомат с состояниями Start, Ident, Int, Real, String, Op, Comment.
Переходы определяются текущим символом; для идентификаторов после
накопления имени выполняется поиск в таблице ключевых слов (регистр
не важен). Ошибки (незакрытая строка, недопустимый символ)
сопровождаются позицией <<строка:столбец>>.

\subsection{Синтаксический анализатор}

Синтаксический анализатор построен методом рекурсивного спуска.
Каждое правило грамматики соответствует функции; в процессе разбора
вызываются семантические процедуры, генерирующие фрагменты ПОЛИЗа.
Ошибки синтаксиса и семантики (необъявленная переменная,
несовпадение типов) выдаются с указанием места.

\subsection{Таблица идентификаторов}

Класс \texttt{SymbolTable} реализует хеш-таблицу с разрешением
коллизий методом цепочек. Хеш-функция — djb2:
\begin{equation}
h(s) = \Bigl(\bigl((5381 \cdot 33) + s_0\bigr)\cdot 33 + s_1\Bigr)\cdots \bmod m,
\end{equation}
где $m = 211$ — простое число. Коллизии разрешаются списками
в каждой ячейке. Средняя сложность поиска — $O(1+\alpha)$, где
$\alpha = n/m$ — коэффициент заполнения.

\subsection{ПОЛИЗ и стековая машина}

Базовый набор операций: JMP, JZ, SET, ADD, SUB, AND, OR, CMPE, CMPL.
Расширения: MUL, DIV, CMPG, CMPLE, CMPGE, CMPNE, NOT, IN, OUT, OUTSTR,
CALL, RET, PUSH\_CONST, PUSH\_VAR, POP\_VAR. Интерпретатор хранит стек
значений, счётчик команд и таблицу идентификаторов; поддерживает
пошаговое выполнение и точки останова.

\subsection{Отладчик}

Реализованы пошаговое выполнение, точки останова, просмотр текущей
инструкции ПОЛИЗа, стека и переменных, редактирование значений
переменных, синхронизация: каждая инструкция ПОЛИЗа хранит ссылку
на (line, col) исходного текста.

% ФОТО: главное окно GUI. Положите в zapiska/fig/gui_main.png
\begin{figure}[H]
  \centering
  \fbox{\parbox[c][8cm][c]{0.9\textwidth}{\centering
    \textbf{[МЕСТО ДЛЯ ФОТО]}\\[4pt]
    Главное окно интерпретатора}}
  \caption{Главное окно программы}
  \label{fig:gui}
\end{figure}

% ФОТО: точка останова. Положите в zapiska/fig/breakpoint.png
\begin{figure}[H]
  \centering
  \fbox{\parbox[c][7cm][c]{0.9\textwidth}{\centering
    \textbf{[МЕСТО ДЛЯ ФОТО]}\\[4pt]
    Остановка на точке останова}}
  \caption{Точка останова и синхронизация с исходным текстом}
  \label{fig:breakpoint}
\end{figure}

% ===================== 4 ТЕСТИРОВАНИЕ =====================
\section{Тестирование}

Набор тест-кейсов приведён в таблице~\ref{tab:tests}.

\begin{table}[H]
\caption{Набор тест-кейсов}
\label{tab:tests}
\centering
\begin{tabular}{@{}clp{7cm}@{}}
\toprule
№ & Вход & Ожидаемый результат \\
\midrule
1 & \texttt{DIM x AS INTEGER; x := 5; PRINT x} & вывод \texttt{5} \\
2 & \texttt{x := 1.5} при \texttt{x AS INTEGER} & ошибка типизации с позицией \\
3 & \texttt{DIM s AS STRING; s := "hi"; PRINT s} & вывод \texttt{hi} \\
4 & \texttt{DO ... LOOP WHILE i < 3} & тело цикла выполняется $\ge 1$ раза \\
5 & деление на ноль & ошибка выполнения \\
6 & \texttt{GOTO undefined} & ошибка парсера \\
7 & незакрытая строка \texttt{"abc} & ошибка лексера \\
8 & \texttt{IF ... THEN ... ELSE ... ENDIF} & корректное ветвление \\
9 & \texttt{SELECT CASE} с диапазоном \texttt{1 TO 5} & корректный выбор ветви \\
10 & \texttt{FUNCTION f(x AS INTEGER) ... ENDFUNCTION} & корректный вызов и возврат \\
\bottomrule
\end{tabular}
\end{table}

% ФОТО: сообщение об ошибке. Положите в zapiska/fig/error.png
\begin{figure}[H]
  \centering
  \fbox{\parbox[c][6cm][c]{0.85\textwidth}{\centering
    \textbf{[МЕСТО ДЛЯ ФОТО]}\\[4pt]
    Сообщение об ошибке с указанием позиции}}
  \caption{Диагностика ошибки}
  \label{fig:error}
\end{figure}

% ===================== ЗАКЛЮЧЕНИЕ =====================
\structtitle{Заключение}

В ходе работы разработан интерпретатор BASIC-подобного языка
(вариант 3-2-1-A) со строгой типизацией, циклом с постусловием,
таблицей идентификаторов на хеш-таблице с цепочками. Реализованы
лексер на конечном автомате, рекурсивный спуск с генерацией ПОЛИЗа,
стековая машина, отладчик и графический интерфейс. Проект
соответствует требованиям на оценку «5».

Возможные расширения: дополнительные типы данных (вектор, список,
ассоциативный массив), вынос стековой машины в отдельный процесс,
защита от копирования.

% ===================== СПИСОК ИСПОЛЬЗОВАННЫХ ИСТОЧНИКОВ =====================
\structtitle{Список использованных источников}

\begin{enumerate}
  \item Ахо А., Лам М., Сети Р., Ульман Дж. Компиляторы: принципы,
        технологии и инструменты. — М.: Вильямс, 2008. — 1184 с.
  \item Вирт Н. Построение компиляторов. — М.: ДМК Пресс, 2010. — 192 с.
  \item Свердлов С. З. Языки программирования и методы трансляции. —
        СПб.: Питер, 2007. — 638 с.
\end{enumerate}

% ===================== ПРИЛОЖЕНИЯ =====================
\structtitle{Приложение А}

\begin{center}
Листинги ключевых модулей
\end{center}

\subsection*{А.1 Модуль lexer.py}
\begin{lstlisting}[language=Python]
class Lexer:
    def tokenize(self):
        while self.pos < len(self.src):
            ch = self._peek()
            if ch in ' \t\r': self._advance(); continue
            if ch == "'": self._comment(); continue
            if ch == '\n': self._eol(); continue
            if ch.isalpha() or ch == '_': self._ident(); continue
            if ch.isdigit(): self._number(); continue
            if ch == '"': self._string(); continue
            self._operator()
        self.tokens.append(Token(TokenType.EOF, None, self.line, self.col))
        return self.tokens
\end{lstlisting}

\subsection*{А.2 Модуль symbol\_table.py}
\begin{lstlisting}[language=Python]
class SymbolTable:
    def __init__(self, size=211):
        self.size = size
        self.buckets = [[] for _ in range(size)]
    def _hash(self, key):
        h = 5381
        for ch in key.lower():
            h = ((h << 5) + h) + ord(ch)
        return h % self.size
    def add(self, sym):
        idx = self._hash(sym.name)
        for s in self.buckets[idx]:
            if s.name.lower() == sym.name.lower():
                raise ValueError("duplicate identifier")
        self.buckets[idx].append(sym)
    def lookup(self, name):
        idx = self._hash(name)
        for s in self.buckets[idx]:
            if s.name.lower() == name.lower():
                return s
        return None
\end{lstlisting}

\subsection*{А.3 Модуль interpreter.py (фрагмент)}
\begin{lstlisting}[language=Python]
def step(self):
    ins = self.poliz[self.pc]
    op = ins[0]
    if op == Op.ADD:
        b = self._pop(); a = self._pop(); self._push(a + b)
    elif op == Op.JZ:
        v = self._pop()
        self.pc = ins[1] if not v else self.pc + 1
        return self.finished
    self.pc += 1
    return self.finished
\end{lstlisting}

\end{document}